\noindent\textbf{Resultado.} Las condiciones físicas son
\[
k_0u_x(0,t)=h_0u(0,t),\qquad
-2k_0u_x(L,t)=h_Lu(L,t).
\]
El problema espacial queda
\[
-\left[k_0\left(1+\frac{x}{L}\right)X'\right]'=\lambda c_0X,
\qquad k_0X'(0)=h_0X(0),\quad
2k_0X'(L)=-h_LX(L).
\]
Por tanto, \(p(x)=k_0(1+x/L)\), \(r(x)=0\), \(w(x)=c_0\). El problema es regular en \([0,L]\): \(p\) es continuamente derivable y positivo, y \(r,w\) son continuos con \(w>0\). El parámetro \(\lambda\) tiene unidades de tiempo inverso.

\noindent\textit{Desarrollo.} El flujo \(q\) se mide respecto de la dirección positiva del eje, que apunta a la derecha. En el extremo izquierdo, la dirección saliente es la contraria: el flujo saliente es \(-q(0,t)\). En el derecho coincide con el eje: el flujo saliente es \(q(L,t)\). Las leyes de transferencia dan
\[
-q(0,t)=h_0u(0,t),\qquad q(L,t)=h_Lu(L,t).
\]
Al usar \(q=-ku_x\), \(k(0)=k_0\) y \(k(L)=2k_0\), resultan
\[
k_0u_x(0,t)=h_0u(0,t),\qquad
-2k_0u_x(L,t)=h_Lu(L,t).
\]
Los signos de las derivadas son distintos porque lo son las direcciones salientes. Las relaciones siguen siendo válidas para un exceso negativo: el flujo saliente firmado puede entonces ser negativo.

El balance interior y la ley de flujo proporcionan
\[
c_0u_t=(ku_x)_x.
\]
Sustituyendo \(u=Xe^{-\lambda t}\) y dividiendo por el factor exponencial,
\[
-\lambda c_0X=(kX')',
\]
es decir,
\[
-(pX')'+rX=\lambda wX,
\qquad p=k_0(1+x/L),\quad r=0,\quad w=c_0.
\]
Las condiciones espaciales son homogéneas y separadas:
\[
h_0X(0)-k_0X'(0)=0,\qquad
h_LX(L)+2k_0X'(L)=0.
\]
En cada extremo los dos coeficientes no son simultáneamente nulos. En el intervalo finito cerrado, \(p\) es continuamente derivable y satisface \(k_0\le p\le2k_0\); \(w=c_0>0\) y \(r=0\) son continuos. Por tanto, se cumplen las condiciones de regularidad. La positividad exigida es la de \(p\) y \(w\); no hace falta que \(r\) sea estrictamente positivo.

Como \(\lambda t\) debe ser adimensional, \(\lambda\) tiene unidades de tiempo inverso. También se obtiene del balance: \(kX''\) representa potencia por unidad de volumen, mientras \(c_0X\) representa energía por unidad de volumen. La identificación no requiere calcular autovalores.
